% pdfLaTeX 编译两次才能解决书签中文乱码
% !TEX TS-program = pdfLaTeX+MakeIndex+BibTeX
% !TEX encoding = UTF-8

\documentclass[CJKutf8,xcolor=x11names]{beamer}

\usecolortheme{seahorse}
\useinnertheme[shadow]{rounded}  
\useoutertheme[subsection=false]{smoothbars}

\usepackage{CJKutf8}
\usepackage{times} % Times New Roman
\graphicspath{{figures/}} % 图片路径
\usepackage{ragged2e}
\renewcommand{\raggedright}{\leftskip=0pt \rightskip=0pt plus 0cm} % 对齐
\usepackage{calligra} % Thank you
\usepackage{fontawesome} % 图标
\usepackage{natbib} % 参考文献

\renewcommand\contentsname{目录}
\renewcommand\listfigurename{插图列表}
\renewcommand\listtablename{表格列表}
\renewcommand\refname{参考文献}
\renewcommand\indexname{索引}
\renewcommand\figurename{图}
\renewcommand\tablename{表}
\renewcommand\partname{部分}
\renewcommand\appendixname{附录}
\renewcommand\abstractname{摘要}

\title{空间广义线性混合效应模型及其应用}

\subtitle{毕业答辩}
\author{导师：李再兴\\学生：黄湘云}
\institute{\includegraphics[width=10ex,interpolate=true]{logo} \\~~ \\ 
理学院\\ 计算数学与统计系 \\ 中国矿业大学（北京）}
\date{ 2018年6月}

\keywords{广义线性模型，空间过程}
\subject{毕业答辩}


\begin{document}
\begin{CJK*}{UTF8}{gkai}  

\maketitle



\begin{frame}

\textbf{Outline 目录}
\begin{enumerate}
  \item Introduction (Motivations and goals)
  \item Literature reviews  
  \item Geostatistical model (SGLMM)
  \item Computing details and simulations
  \item Real data analysis (Applications)
  \item Discussion
\end{enumerate}
\end{frame}

\section{介绍}

\begin{frame}
\textbf{Outline 目录}
\begin{enumerate}
  \item Introduction (Motivations and goals)
  \item Literature reviews  
  \item Geostatistical model (SGLMM)
  \item Computing details and simulations
  \item Real data analysis (Applications)
  \item Discussion
\end{enumerate}
\end{frame}

\section{$\chi^2$极限分布}

\begin{frame}
\textbf{Outline }
\begin{enumerate}
  \item Introduction (Motivations and goals)
  \item Literature reviews  
  \item Geostatistical model (SGLMM)
  \item Computing details and simulations
  \item Real data analysis (Applications)
  \item Discussion
\end{enumerate}
\end{frame}




\begin{frame}[allowframebreaks]
\frametitle{Frameworks, Packages and Softwares}

\begin{description}
\item[\textbf{R:}] geoR geoRglm spatial PrevMap \\
\citet{R-geoR,R-geoRglm,R-spatial,Giorgi2016}

\item[\textbf{Stan:}] Stan \footnote{\scriptsize \url{http://mc-stan.org/} }  
interfaces with R (RStan) ,Python (PyStan) , MATLAB (MatlabStan) 
and  more \\ \citet{Stan2015,Stan2017}

\item[\textbf{PyMC3:}] Probabilistic programming in Python using PyMC3 \\ \citet{Salvatier2016}

\item[\textbf{JAGS:}] \textbf{J}ust \textbf{A}nother \textbf{G}ibbs \textbf{S}ampler \footnote{\scriptsize \url{https://en.wikipedia.org/wiki/Just_another_Gibbs_sampler}} \\
Bayesian hierarchical models using Markov chain Monte Carlo (MCMC)
\item[\textbf{BUGS:}] \textbf{B}ayesian inference \textbf{U}sing \textbf{G}ibbs \textbf{S}ampling ,
such as winBUGS, OpenBUGS 
\item[\textbf{R-INLA:}] \textbf{I}ntegrated \textbf{N}ested \textbf{L}aplace \textbf{A}pproximations \\
\citet{Rue2009,R-INLA,Rue2017arXiv}
\end{description}


\end{frame}

\section*{致谢}

\begin{frame}
\centering
\color{red}
\Huge {\calligra{Thank You}}\\
\bigskip
\color{Blue4}
\normalsize
\faEnvelope{} \href{mailto:me@somewhere.com}{xiangyunfaith@outlook.com} \\
\faGithub{} \href{https://github.com/Cloud2016}{https://github.com/Cloud2016} \\
\bigskip
\includegraphics[width=10ex,interpolate=true]{me}
\end{frame}


\begin{frame}[allowframebreaks]
\bibliographystyle{authordate1}
\bibliography{R-GLMM-pkgs}
\end{frame}

\end{CJK*}
\end{document} 


